% ===================================================================
\chapter{Hardware Implementation}
\label{ch:hardware}

Both algorithms were implemented as a compressor core and a matching
decompressor core in C++ for Vitis HLS, and built into a test system on an AMD
KCU116 evaluation board. This chapter describes the cores, how they stay
exactly lossless on any input, and the board system used to run and measure
them.

\section{Common interface and bit packing}
\label{sec:interface}

All four cores share one interface. A compressor takes in 128 bits per clock,
one complex 16-bit sample from each of four receive channels, and sends out the
compressed stream in 256-bit words, with the last word of each frame marked; a
decompressor does the reverse. The frame dimensions are set through a control
interface, so the compressed stream needs no header.

Both compressors share one bit packer. It collects codewords in a 512-bit
register and writes out a 256-bit word whenever one is full, so only the last
word of each frame is padded. If a frame happens to end exactly on a word
boundary, an empty word is sent to mark the end. The codewords are stored
bit-reversed, because the packer fills from the least significant bit; the
reversed code is still prefix-free. The Huffman dictionary is a sixteen-entry
table, so each core was built with both dictionaries of
Section~\ref{sec:dict}, which differ only in the contents of that table.

\section{DRHE compressor}
\label{sec:drheimpl}

DRHE keeps five floating-point state values for each transmitter, receive
channel and range bin: the magnitude and phase predictions, the previous
magnitude and the previous two phases. The transmitter index is the ramp number
modulo $\nTx$, so predicting from the same transmitter
(Section~\ref{sec:predaxis}) only changes which memory address is read. Each
transmitter's state is reset on its first ramp of every frame, so every frame
decodes on its own. The four receive channels are four parallel copies of the
datapath, and the state fills 40 block RAMs.

For every sample, DRHE converts the reconstructed value to polar form with
\texttt{sqrt} and \texttt{atan2} and converts the prediction back with
\texttt{sin} and \texttt{cos}, in single-precision floating point. These four
operators make up about two thirds of the logic and all 132 DSP slices;
Section~\ref{sec:choice} discusses this choice against Kiem's fixed point.

The core accepts a new input every two clocks ($\mathrm{II} = 2$). Each state
memory has two ports, and on the first ramps of a frame the loop needs a third
access to reset the state, so the tool schedules every iteration over two
clocks. Removing that extra write reaches one clock per input, but doubles the
floating-point hardware, and the design is already far faster than the radar
needs (Section~\ref{sec:boardtiming}).

\section{LPC compressor}
\label{sec:lpcimpl}

LPC's coefficients depend on the whole frame, so its compressor runs in three
passes (Figure~\ref{fig:lpcblock}). It first stores the frame in a buffer.
Then, for each receive channel and range bin, it computes the autocorrelations
at lags $0$, $\nTx$ and $2\nTx$, solves Equation~\eqref{eq:yw}, and overwrites
the buffer with the residuals, so one \SI{3.0}{\mebi bit} memory serves both
purposes. Finally it sends the coefficient table and then the coded residuals.
Every loop accepts one value per clock. A frame takes \num{325366} clocks, of
which about a fifth is the pipeline filling at the start of each of the 512
per-bin loops.

\begin{figure}[t]
\centering
\begin{tikzpicture}[font=\small,
  b/.style={draw=slate,fill=white,rounded corners=2pt,minimum width=30mm,minimum height=9mm,font=\scriptsize,align=center},
  h/.style={b,draw=accent,line width=0.9pt,fill=lightbg},
  m/.style={draw=slate,fill=warm!12,rounded corners=2pt,minimum width=18mm,minimum height=22mm,font=\scriptsize,align=center},
  a/.style={-{Latex[length=1.5mm]},slate}]
\node[b]  (in)  at (0,0)    {AXI in\\128\,b/beat};
\node[h]  (p1)  at (3.3,0)  {\textbf{pass 1}\\ingest};
\node[m]  (buf) at (6.6,0)  {frame\\buffer\\\SI{3.0}{\mebi bit}};
\node[h]  (p2)  at (9.9,0)  {\textbf{pass 2}, per bin\\correlate, solve,\\residual (in place)};
\node[h,minimum width=44mm]  (p3)  at (6.6,-2.5) {\textbf{pass 3}: emit coefficients,\\then residuals in raster order};
\node[b]  (out) at (0,-2.4) {AXI out\\256\,b/beat};
\draw[a] (in) -- (p1);
\draw[a] (p1) -- (buf);
\draw[a] (buf) -- (p2);
\draw[a] (p2.south) -- (9.9,-2.5) -- (p3.east);
\draw[a] (p3) -- (out);
\node[font=\scriptsize,text=slate,align=center] at (9.9,1.15) {$R_0,R_1,R_2$ at lags $0,\nTx,2\nTx$};
\end{tikzpicture}
\caption[LPC three-pass architecture]{The LPC compressor. The Yule--Walker
coefficients depend on the whole frame, so the compressor buffers it. Each
(channel, bin) is then correlated, solved and turned into residuals, which
overwrite the buffer in place so that only one memory is needed. Emitting the
residuals in the original raster order lets the decompressor work without a
frame buffer.}
\label{fig:lpcblock}
\end{figure}

The coefficients are quantised to 16-bit fixed point ($a_1$ in $[-2, 2)$ and
$a_2$ in $[-1, 1)$, following Equation~\eqref{eq:stability}), which changes no
residual over the 50 frames. The autocorrelations are summed as exact integers,
because when the data is strongly correlated the solve subtracts two nearly
equal quantities and floating point would lose precision.

The residuals are sent in the order the samples arrived, not the order the
second pass produced them. The decoder therefore needs no frame buffer: it
keeps the coefficients and the last $2\nTx$ samples of each channel and bin,
and outputs each sample as its residual arrives. LPC's need for the whole frame
costs only the compressor.

\section{Exactly lossless on any input}
\label{sec:robust}

Real radar data never reaches the top of the coder's range
(Section~\ref{sec:edgecases}), but synthetic input can, and one value needs
care. Category 15 holds exactly the $2^{15}$ payload patterns of
$|v| \in [16384, 32767]$, so a residual of $-32768$, which 16-bit arithmetic can
produce, has no code of its own and would be confused with $+32767$. The
compressors saturate it to $-32767$, which costs one least significant bit on
that sample, keeps the format unchanged, and never happens on real data.

Saturation alone would make the encoder and decoder disagree, because the
decoder rebuilds the saturated value and DRHE's recursive filters would carry
the difference through every later ramp. The DRHE encoder therefore rebuilds
each sample exactly as the decoder will and updates its predictor from that
rather than from the original. On real data this changes no output bit; its
cost is 26 of the pipeline's 58 stages.

\section{Decompressors}
\label{sec:dechw}

A decompressor was built for each algorithm, so that every frame can be checked
bit for bit and the whole sensor-to-processor pipeline exists in hardware. Each
is also compiled as a PC program from the same C++ source. Decoding is serial,
since a codeword's length is only known once it has been matched, so a decoder
that kept pace with the compressor would need eight chained matches per clock
and, for DRHE, four copies of the floating-point predictor. Both decoders were
instead designed to be small: they decode one value per clock, about
\SI{2}{\milli\second} per frame, which is still a hundred times faster than the
radar produces frames.

Each decoder has two stages connected by a stream. The first compares the input
against all sixteen codewords at once and extracts one value per clock. The
second rebuilds the samples. In the DRHE decoder this is one floating-point
predictor using exactly the compressor's expressions, since a recursive
predictor drifts if the two sides disagree even in the last bit; a
double-precision Python decoder written as a cross-check drifts on 8 of the 50
frames. In the LPC decoder it is two fixed-point multiply--accumulates per
part, with the past samples in distributed RAM, so it needs no block RAM.
After routing the DRHE decoder uses \num{14155} LUTs and the LPC decoder
\num{15212}.

\section{Range FFT on chip}
\label{sec:fehw}

The compressors take FX16 range-FFT codes, which the method study takes from
MATLAB (Section~\ref{sec:refchain}). To check that no result depends on this,
and to let the board start from raw ADC samples as a sensor would, the MATLAB
range FFT was also written in HLS. For each ramp of each of the four channels
it removes the ramp mean, applies the Hann window and a 256-point fixed-point
FFT, keeps bins 0 to 127 and rounds them to FX16 as Equation~\eqref{eq:scale}
does. It sits between the input DMA and the compressor, takes about 1.5 clocks
per input word, and uses \num{14511} LUTs and 296 DSP slices. Fixed-point
arithmetic cannot match double precision exactly, so a value very close to a
rounding boundary can round the other way; Section~\ref{sec:feresults} measures
how often.

\section{The board test system}
\label{sec:system}

The cores run in a complete system on the KCU116 board
(Figure~\ref{fig:system}). A MicroBlaze processor controls it, and \SI{1}{\giga
\byte} of DDR4 memory holds the input, compressed and reconstructed frames. One
DMA streams each frame from memory through the compressor and writes the
compressed stream back. In the \texttt{\_full} builds a second DMA streams it
through the decompressor, so the board compresses, decompresses and checks
every frame itself, and the \texttt{\_fe\_full} builds add the range FFT. For
each frame the firmware starts the transfers and reads a hardware timer at the
start, at the last input word and at the last output word, which gives the time
per frame and the latency of Section~\ref{sec:metrics}. A script on the PC
loads the frames, runs the firmware, reads back the results, and checks every
compressed frame byte for byte against the C-simulation and every
reconstruction against the original.

\begin{figure}[H]
\centering
\begin{tikzpicture}[font=\small,x=0.88cm,
  b/.style={draw=slate,fill=white,rounded corners=2pt,minimum height=8mm,inner xsep=4pt,font=\scriptsize,align=center},
  h/.style={b,draw=accent,line width=0.9pt,fill=lightbg},
  a/.style={{Latex[length=1.4mm]}-{Latex[length=1.4mm]},slate}]
\node[b,fill=lightbg,minimum width=20mm] (pc) at (-3.4,1.7) {host PC\\(JTAG, UART)};
\node[b,minimum width=24mm] (mb) at (0,1.7) {MicroBlaze\\128\,KB local RAM};
\node[b,minimum width=24mm] (io) at (0,0.1) {UART, timer,\\LEDs};
\node[b,minimum width=20mm,minimum height=32mm] (ic) at (3.4,0.9) {AXI\\SmartConnect};
\node[b,minimum width=24mm] (ddr) at (7.0,2.5) {DDR4-2400, 1\,GB\\frames and results};
\node[b,minimum width=18mm] (dc) at (7.0,0.9) {DMA \texttt{dma\_c}};
\node[b,minimum width=18mm] (dd) at (7.0,-0.7) {DMA \texttt{dma\_d}};
\node[h,minimum width=28mm] (cc) at (10.6,0.9) {compressor\\128\,b in, 256\,b out};
\node[h,minimum width=28mm] (dk) at (10.6,-0.7) {decompressor\\256\,b in, 128\,b out};
\draw[a] (pc) -- (mb);
\draw[a] (mb) -- (mb -| ic.west);
\draw[a] (io) -- (io -| ic.west);
\draw[a] (ic.east |- ddr) -- (ddr);
\draw[a] (ic.east |- dc) -- (dc);
\draw[a] (ic.east |- dd) -- (dd);
\draw[a] (dc) -- (cc);
\draw[a] (dd) -- (dk);
\node[font=\scriptsize,text=slate] at (10.6,-1.45) {\texttt{\_full} builds only};
\end{tikzpicture}
\caption[The KCU116 test system]{The board test system. The host places the
frames in DDR4 over JTAG; the MicroBlaze starts the DMAs, times each frame and
writes a result record; the host then reads the results and compressed frames
back. Everything runs at \SI{100}{\mega\hertz}.}
\label{fig:system}
\end{figure}

\section{Building and timing closure}
\label{sec:timing}

Each core is exported from HLS as an IP block and built into the system with
Vivado, six bitstreams in all. After placement and routing all six meet
\SI{100}{\mega\hertz} with positive slack (Table~\ref{tab:builds}).

\begin{table}[H]
\centering
\caption[The six board builds]{The six board builds, whole system (processor,
DDR4 controller, DMAs and cores), post-route on the \texttt{XCKU5P} at
\SI{100}{\mega\hertz}, with the re-derived dictionaries; the original-dictionary
builds also meet timing. BRAM is in 36\,Kb blocks.}
\label{tab:builds}
\small
\begin{tabular}{l|c|c|c|c}
\textbf{Build} & \textbf{LUT} & \textbf{BRAM} & \textbf{DSP} & \textbf{WNS (ns)}\\
\hline
DRHE, compress only    & \num{61011} (\SI{28}{\percent}) & 104.5 & 138 & $+0.243$\\
LPC, compress only     & \num{59203} (\SI{27}{\percent}) & 162.5 & 164 & $+0.174$\\
DRHE, \texttt{\_full}  & \num{81573} (\SI{38}{\percent}) & 144   & 184 & $+0.341$\\
LPC, \texttt{\_full}   & \num{80986} (\SI{37}{\percent}) & 179.5 & 187 & $+0.220$\\
DRHE, \texttt{\_fe\_full} & \num{101383} (\SI{47}{\percent}) & 211 & 480 & $+0.340$\\
LPC, \texttt{\_fe\_full}  & \num{100709} (\SI{46}{\percent}) & 246.5 & 483 & $+0.374$\\
\end{tabular}
\end{table}
